% SPDX-License-Identifier: Apache-2.0
\section{The Same Design, Three Ways}
\label{sec:worked-example}

A multiply accumulate unit reads pairs of words from memory, accumulates
their products, and answers over a bus. It is small enough to print and
large enough to need a contract, a protocol body, a tag and a binding.

\subsection{As TxHDL states it today}

TxHDL writes the protocol well and cannot state the tag, the contract or the
binding.

\begin{lstlisting}
module MacUnit {
    port host: MacBus.slave;
    port mem:  MemBus.master;

    var acc: u64 = 0;

    pipeline mac(a: u32, b: u32, prev: u64) -> u64 {
        stage mul { var p = zext<u64>(a) * zext<u64>(b); }
        stage add { return prev + p; }
    }

    sequence main {
        loop {
            var req = await host.request;
            acc = 0;
            for (i in 0..req.count) {
                mem.request <- { .addr = req.addr + i*8,
                                 .write = false };
                var lo = await mem.response;
                mem.request <- { .addr = req.addr + i*8 + 4,
                                 .write = false };
                var hi = await mem.response;
                acc = mac(lo.data[31:0], hi.data[31:0], acc);
            }
            host.response <- { .acc = acc };
        }
    }
}
\end{lstlisting}

\subsection{As LHDL states it today}

LHDL states the contract, the tag and the binding, and has no way to write
the loop. The listing below is where an LHDL author stops.

\begin{lstlisting}
interface MacOp {
    addr  : uint<32>;
    count : uint<8>;
    acc   : uint<64>;

    modport Server { in addr; in count; out acc; }
}

tag @MemFetch with handshake, capacity=8;

pipeline Mac implements MacOp.Server {
    pipe p: u64;
    stage Mul { p := a * b; }
    stage Add { acc := p + acc; }
}

config Build {
    bind Top.mac_unit => Mac;
}
\end{lstlisting}

No LHDL construct reads \code{count} words one after another and waits for
each response. A state machine would state it as hand written states, which
is the form \code{await} exists to replace.

\subsection{Merged}

\begin{lstlisting}
transaction MacRequest { addr: u32, count: u8 }
transaction MacResponse { acc: u64 }

bus MacBus {
    channel request:  MacRequest,
    channel response: MacResponse,
}

pub trait MacOp {
    async fn run(&mut self, req: MacRequest) -> MacResponse;
}

// Two stages, so the multiplier does not sit in the adder's
// path. The compiler places the register between them.
pipeline mac(a: u32, b: u32, prev: u64) -> u64 {
    stage mul { let p = (a as u64) * (b as u64); }
    stage add { prev + p }
}

#[tag(MemFetch, handshake, capacity = 8)]
pub module MacUnit {
    port host: MacBus::Slave,
    port mem:  MemBus::Master,

    let mut acc: u64 = 0;

    seq main {
        loop {
            let req = host.request.await;
            acc = 0;

            for i in 0..req.count {
                let lo = mem.read(req.addr + i*8).await;
                let hi = mem.read(req.addr + i*8 + 4).await;
                acc = mac(lo.data[0..32], hi.data[0..32], acc);
            }

            assert!(req.count > 0, "empty accumulate");
            host.response.send(MacResponse { acc }).await;
        }
    }
}

impl MacOp for MacUnit {}

config Build {
    bind top.mac_unit => MacUnit;

    domain MemFetch {
        clock: sys_clk_100mhz,
        reset: sys_rst_n,
    }
}
\end{lstlisting}

Nine differences from the TxHDL version are worth naming, because each is a
decision from Table~\ref{tab:syntax}.

\begin{enumerate}
  \item \code{var acc} becomes \code{let mut acc}, so the text says the
        value changes and not that a register exists.
  \item The zero-extension call becomes an \code{as} cast.
  \item The descending bit slice becomes a range.
  \item The prefix await becomes postfix.
  \item The send becomes a call that returns a value, so a request and its
        response are one line instead of two.
  \item The designated initialiser becomes a struct literal, with the type
        named and field shorthand doing the rest.
  \item \code{sequence} becomes \code{seq}, and its cycle counts stay
        exactly as written.
  \item The tag, the trait and the configuration come from LHDL, and TxHDL
        had no way to write any of them.
  \item The assertion comes from TxHDL, and LHDL had no way to write it.
\end{enumerate}
